\section{A noisy certificate without a reference table}\label{sec:uniform}
The two difficulties in applying \eqref{eq:defect} to noisy data are
identifying body types without a known witness and recovering a discrete
subgroup before its covolume has been calibrated. The first requires
only shape separation, not isolation of every edge key. For the second
we use bounded-denominator separation and postpone the area test until
the subgroup has already been recovered.

\subsection{Uniform physical class}
Fix an upper bound $r_0$ for the number of obstacle types, a diameter
bound $D_0$, and a bound $R$ for every recorded gap. The actual number
and identities of types are unknown. Require $A_0\le A\le A_1$, a
positive lower bound $V_0$ for the true cell covolume, and a positive
lower bound $a_0$ for every obstacle area. One may take $V_0=A_0$;
we keep the symbol separate to display its arithmetic role. There are
uniform positive separation, curvature, channel-isolation and adjacent
free-flight bounds. Lattice bases and their inverses admit uniform bounds
in some reduced presentation. Each observed record satisfies the physical
$C^3$, active, twist, diagonal non-grazing, time-separation and rectangle
conditions of Proposition~\ref{prop:local-stable}, and its densities
have a common $C^{2,\beta}$ bound on a slightly larger rectangle.

After Steiner centering, each support function $p$ is holomorphic and
bounded by $M$ on $|\operatorname{Im}\varphi|<\rho$, for fixed
$M,\rho>0$. Local contact graphs have common bounded holomorphic
extensions on a fixed complex disk. These graph bounds also follow
on smaller disks from the support bounds and positive curvature.
For positive $\eta,\zeta$ require
\begin{align}
 \|p(\cdot+\phi)-p\|_2&\ge
            \eta\dist(\phi,2\pi\Z),\label{eq:rotation}\\
 \|p(-\cdot+\phi)-p\|_2&\ge\eta,\label{eq:reflection}\\
 \inf_{O\in O(2)}\|p_i-p_{OC_j}\|_2&\ge\zeta\quad(i\ne j).
                                                   \label{eq:shape-separation}
\end{align}
Take a closed class with these margins, compact after restricting to a
narrower strip and bounded root placement. Bounds and margins, but no
reference shape, period basis, actual obstacle count or reference witness,
are inputs. The class need not lie in one local deformation chart.
The margins allow infinite-dimensional analytic shape variation.

All comparisons below allow one common Euclidean motion, a permutation
of obstacle types, and a change of representative copies and period basis.
A table error at most $t$ means that such representatives and bases can
be chosen so that all support-function $C^2$ differences and the period
matrix difference are at most $t$. We prove estimates in these explicit
simultaneous gauges; no continuous choice of a canonical lattice basis
across its reduction walls is asserted.

\subsection{From densities to complete pair images}
\begin{lemma}[Uniform pair recovery]\label{lem:pair-error}
If fitted physical density pairs have $C^2$ error at most $e<1$ on
the common rectangles, each recovered entire analytic pair in its own
contact frame has support-function $C^2$ error at most
$\Delta=C e^\theta$, where $C>0$ and $0<\theta<1$ depend only
on the stated class. The recovered free area has error at most $Ce$.
\end{lemma}
\begin{proof}
Proposition~\ref{prop:local-stable} gives $Ce$ errors for the two
intrinsic arcs and $A$. Positive curvature converts a smaller arc to
a normal-angle interval of a fixed positive length. The inverse angle
map and its support values depend locally Lipschitzly on the $C^2$ arc,
using the uniform physical $C^3$ bound. Hence the support functions
are $Ce$-close on a fixed real interval in the contact frame.

Here is the continuation estimate. Slit a narrower periodic strip
cylinder along a shorter observed interval. The harmonic function equal
to one on both slit sides and zero on the two outside boundary circles
is positive in the slit cylinder. On a still narrower closed cylinder,
including the slit by continuity, it has a positive minimum $\theta$.
The endpoints of the slit are regular Dirichlet boundary points.
Apply the subharmonic maximum principle to the logarithm of the modulus
of the support difference, using $Ce$ on the slit and its uniform
complex bound elsewhere. This gives $C'e^\theta$ on the narrower
cylinder. Zeros follow by regularization; Cauchy estimates on fixed
disks give the same bound through two derivatives on the real circle.
Shifting from centered to contact frames changes the complex support
bound by a controlled diameter term. Constants are uniform under rotations
and under translating the observed interval. This is conditional analytic
continuation of functions on arcs, not an infinite jet estimate; compare
\cite{Trefethen}.
\end{proof}

Fitting can be defined independently for each record over a fixed compact
physical pair class. Given a density estimate with certified error $e$,
choose a point of a fixed countable dense family whose $C^2$ discrepancy
is within $e$ of the infimum. It then differs from the true density by
at most $3e$. A first-eligible-point rule is measurable. This is an
existence construction, not an efficient analytic continuation algorithm.
It does not declare every noisy pair of functions to be a physical law.
Uniform physical branch validity remains a model assumption.

\begin{lemma}[Incidence and placements without edge-key neighborhoods]
\label{lem:reference-free-assembly}
For sufficiently small $\Delta$, the recovered body shapes identify
all visible types without a reference table. A spanning tree chosen in
the resulting observed multigraph gives representative and cycle-vector
errors at most $C\Delta$, uniformly over repetitions and surplus edges.
No classification of distinct channel orbits is required.
\end{lemma}
\begin{proof}
Use the centered support distance modulo $O(2)$ in $L^2$. Two estimates
of the same body type are at distance at most $C\Delta$; estimates
of different true types are at distance at least $\zeta-C\Delta$.
For $C\Delta<\zeta/4$, joining observations whose distance is less
than $\zeta/2$ gives exactly the true type clusters. There is no
ambiguous chain of clusters, since the two cases have disjoint distance
ranges. Each record is retained as one edge with its two recovered
endpoint types. Multiple edges and loops are allowed, so near-coincidence
of two different edge geometries never has to be resolved.

For a tree placement, choose a least-discrepancy orthogonal match of
centered source shapes and then match their Steiner points. Compactness
of $O(2)$ ensures existence. A true matching has discrepancy $O(\Delta)$,
so the minimizing match does too. Relative to the true matching its
orthogonal action moves the true support by $O(\Delta)$ in $L^2$.
Equation~\eqref{eq:reflection} excludes the wrong component, and
\eqref{eq:rotation} bounds the remaining rotation angle by
$O(\Delta/\eta)$. The uniform $C^3$ support bound controls the
resulting $C^2$ placement error. Steiner-point translations are bounded
linear functionals of support values. Induction along a tree path of
at most $r_0-1$ edges proves the representative bound.

For every non-tree edge, match its source in the same way and subtract
the Steiner point of the placed target. The corresponding exact physical
operation gives $d_e\in\Lambda$, and the displacement error is
$O(\Delta)$. Constants depend on $r_0$ and the margins, not the
number of repeated or unused records. This comparison uses the truth
only in the proof; the choice of root, tree and matches is entirely
from the fitted observations.
\end{proof}

\subsection{Rational reconstruction before calibration}
Set
\begin{equation}\label{eq:denominator-bound}
 B_0=(2r_0-1)(R+2D_0),\qquad
                  Q=\max\{1,\lceil B_0^2/V_0\rceil\}.
\end{equation}
Every exact non-tree displacement has norm at most $B_0$. Indeed
representative Steiner points in a rooted tree lie within
$(r_0-1)(R+2D_0)$ of the root; an edge pair has Steiner-point distance
at most $R+2D_0$. Subtracting the placed target gives the stated bound.

\begin{lemma}[Reference-free discrete recovery]\label{lem:rational}
Suppose estimated displacements have uniform error at most $t$ and
$t$ is sufficiently small in terms of $B_0,V_0$. One can determine
whether their true rank is two. In the rank-two case choose a pair
$D$ by maximizing the estimated absolute determinant. Each coordinate
of $D^{-1}d_e$ is a rational with denominator at most $Q$. For all
sufficiently small $t$, these coordinates, their least common denominator
$q$, and the integer column group in \eqref{eq:denominator-group} are
recovered exactly by a finite rational search, using no area calibration.
Moreover $q\le Q$, and the estimated cycle covolume has error at most
$Ct$.
\end{lemma}
\begin{proof}
The determinant error of each pair is at most $(2B_0+t)t$. In a true
rank-two component, every nonzero pair determinant is an integer multiple
of $\covol\Lambda$ and hence at least $V_0$. In a rank-one component
all pair determinants vanish. Once the error is less than $V_0/4$,
a maximal estimated determinant exceeding $V_0/2$ selects an independent
true pair, and no dependent component passes this test. The selected
true determinant lies between $V_0$ and $B_0^2$.

Write $n_D=|\det D|/\covol\Lambda\in\{1,\ldots,Q\}$. The
finite-group argument preceding \eqref{eq:denominator-group} puts each
coordinate in $n_D^{-1}\Z$. Its absolute value is bounded, for example,
by $2B_0^2/V_0$. Matrix inversion on this determinant-separated set
gives coordinate error at most $C_0t$. Distinct reduced rationals with
denominators at most $Q$ differ by at least $Q^{-2}$: their difference
has a nonzero integer numerator over a denominator at most $Q^2$.
If $C_0t<1/(4Q^2)$, the true value is the unique rational in the
allowed bounded range within $1/(4Q^2)$ of the estimate. A finite
search over denominators and bounded numerators recovers it exactly.
All actual denominators divide the same $n_D$, so their least common
multiple $q$ divides $n_D$ and is at most $Q$.

Hermite reduction now gives the same integer column group and a basis
$H$ for the truth and its estimate. Since this group contains
$q\Z^2$, it has $1\le|\det H|\le q^2$. Formula
\eqref{eq:denominator-group} with the estimated $D$ gives the covolume
estimate; its error is at most the pair-determinant error. All computations
of integer structure precede any use of $A+S_I$.
\end{proof}

This is not integer spanning of noisy real vectors, which could give
a dense subgroup. The finite denominator bound comes from periodicity
and a covolume lower bound; rational separation fixes the exact arithmetic
first. The search may be large, but is finite once the stated priors
are fixed. It does not make the preceding physical fitting efficient.

\subsection{The decision and reconstruction}
After type clustering, consider each connected component. Reject rank
less than two. In rank two compute $q,H$ as above, and put
\begin{equation}\label{eq:estimated-defect}
 \widehat{\mathcal D}
 =\frac{|\det\widehat D|\,|\det H|}{q^2}
           -\widehat A-\sum_{i\in I}\area(\widehat C_i).
\end{equation}
The estimate $\widehat A$ may be taken from any one fitted record;
all such estimates are uniformly close to the common true value. Use
one fitted placed body per type, rather than summing repetitions.

\begin{proof}[Proof of Theorem~\ref{thm:uniform-main}]
Lemmas~\ref{lem:pair-error}--\ref{lem:rational} give the correct visible
clusters, graph and rational cycle structure when $e<e_0$ for a
class-dependent threshold. For a $C^2$ support function,
\[
 \area(C)=\tfrac12\int_0^{2\pi}(p^2-(p')^2)\dd\varphi.
\]
On the bounded class this is Lipschitz in the $C^1$ norm. There are at
most $r_0$ visible types, so \eqref{eq:estimated-defect} differs from
\eqref{eq:defect} by at most $C\Delta$. Reduce $e_0$ until this is
less than $g_0/4$, with $g_0=\min(V_0,a_0)$. Accept precisely when
\begin{equation}\label{eq:acceptance}
                     |\widehat{\mathcal D}|<g_0/2.
\end{equation}
Corollary~\ref{cor:defect-gap} shows that every complete saturated
component passes and every incomplete component fails. Neither witness
retention nor coverage is assumed in proving soundness of the decision.

On acceptance all types are visible and $\Gamma=\Lambda$. The
reconstructed period basis is $\widehat B=\widehat D H/q$, while
$B=DH/q$ is a basis of the true period lattice. The integer group
contains $q\Z^2$; there are only finitely many such column Hermite
forms for $q\le Q$, so $\|H/q\|$ has a bound depending only on $Q$.
Thus $\|\widehat B-B\|\le C_Q\Delta$. The placed body supports
are simultaneously $C\Delta$-close by Lemma~\ref{lem:reference-free-assembly}.
This gives the stated whole-table bound in common representative and
period gauges. Small errors and the strict curvature and global separation
margins keep this finite representation a physical periodic table:
bounded bases and their inverses reduce potentially close translated
pairs to a uniformly finite set; the remaining pairs are farther apart.
Alternatively one can project to the compact admissible physical class
with an additional factor in the error, since the true table is within
that error of the construction. A countable dense near-minimizer gives
a measurable selection. All decisions use the observations and global
bounds, never a reference table or local witness-key neighborhoods.
\end{proof}

The decision is uniform over any data-dependent deletion of records on
a simultaneous accuracy event for the original finite list. Surplus
channels need no common key separation and may change as long as the
ones actually used remain genuine branches within the local physical
class. No theorem is asserted for fabricated records or an unrestricted
mixture of unseparated itineraries. A failed test is an informative
noncertificate, not evidence that a table does not exist.
